\documentclass[11pt]{article}

\usepackage[a4paper,margin=1in]{geometry}
\usepackage{amsmath,amssymb,amsthm,mathtools}
\usepackage{booktabs}
\usepackage{enumitem}
\usepackage{microtype}
\usepackage{xcolor}
\usepackage{hyperref}

\hypersetup{
    colorlinks=true,
    linkcolor=blue!60!black,
    citecolor=blue!60!black,
    urlcolor=blue!60!black,
    pdftitle={Quotients Without Division},
    pdfauthor={Vinicius Santos}
}

\newtheorem{definition}{Definition}[section]
\newtheorem{theorem}[definition]{Theorem}
\newtheorem{proposition}[definition]{Proposition}
\newtheorem{corollary}[definition]{Corollary}
\newtheorem{lemma}[definition]{Lemma}
\newtheorem{remark}[definition]{Remark}
\newtheorem{question}[definition]{Question}
\newtheorem{conjecture}[definition]{Conjecture}

\newcommand{\B}{\mathsf{B}}
\newcommand{\Tr}{\operatorname{Tr}}
\newcommand{\Nm}{\operatorname{N}}
\newcommand{\Sp}{\operatorname{Sp}}
\newcommand{\Quot}{\operatorname{Quot}}
\newcommand{\Fact}{\operatorname{Fact}}
\newcommand{\Const}{\operatorname{Const}}
\newcommand{\ph}{\varphi}
\newcommand{\ps}{\psi}
\newcommand{\q}{q}
\newcommand{\Z}{\mathbb{Z}}
\newcommand{\Q}{\mathbb{Q}}
\newcommand{\R}{\mathbb{R}}
\newcommand{\dd}{\Delta_\ph}
\newcommand{\pd}{\partial_\ph}
\newcommand{\Cphi}{C_\ph}

\title{Quotients Without Division\\
\large Unit Actions, Relations, and Factor Projections in a $\varphi$-Native Algebra}
\author{Vinicius Santos\\
\small Exploratory draft}
\date{June 2026}

\begin{document}
\maketitle

\begin{abstract}
The companion manuscript \emph{A Golden Shadow Algebra} isolates a division-free scalar grammar generated by the golden ratio $\ph$ and the operation
\[
    \B(x,y)=xy-y.
\]
Its exact constants are $\Z[\ph]$, countable and dense in a chosen real embedding of $\R$ but discrete under the full Minkowski embedding, and its scalar expressions are polynomial and total.  Thus the $\ph$-native system achieves density without introducing a division operation.  This note asks what remains of division once density no longer requires it.

The answer proposed here is that division is not one operation.  It is an overloaded bundle of roles: unit scaling, fixed-denominator rates, rational constants, reciprocal functions, quotient solving, derivative quotients, and normalization.  In a $\ph$-native algebra these roles split.  Division by $\ph$ is not a singular reciprocal but a unit action,
\[
    x/\ph=(\ph-1)x=\B(\ph,x).
\]
Exact non-unit reciprocals are not finite scalar terms; they may be carried as rate classes, evaluated in controlled localizations, or approximated by dense $\Z[\ph]$-constants.  Reciprocal functions become local polynomial approximation problems.  Quotients become relations $u=qv$.  Derivatives become factor projections: a golden difference is divided not by a primitive quotient, but by extracting the factor of the golden displacement.

The central worked example is the golden contraction
\[
    \Cphi(x)=\B(\ph,x)=(\ph-1)x.
\]
Writing $q=\ph^{-1}=\ph-1$, the golden difference
\[
    \dd f(x)=f(x)-f(qx)
\]
is divisible by $\dd x=x-qx$ for every polynomial $f\in\Z[\ph][x]$.  The derivative-like operator $\pd f$ is defined as the unique factor
\[
    \dd f=(\pd f)\,\dd x,
\]
not as a primitive quotient.  For monomials this gives
\[
    \pd x^n=[n]_\ph x^{n-1},
    \qquad
    [n]_\ph=1+q+\cdots+q^{n-1}\in\Z[\ph].
\]
This overlaps formally with Jackson's $q$-calculus at $q=\ph^{-1}$, but the point of the paper is different: the $q$-difference structure appears as a consequence of a division-free, leaf-generated grammar and its factor-projection mechanism.
\end{abstract}

\tableofcontents

\section{Introduction: division is not one operation}

In ordinary arithmetic, division is written as a single symbol:
\[
    \frac{u}{v}.
\]
Yet this symbol performs several different jobs.  It scales by units, introduces rational constants, defines reciprocal functions, solves multiplicative equations, normalizes objects, and forms slopes.  Treating all of these jobs as one primitive operation is convenient, but it also imports denominator singularities.  Expressions such as
\[
    \frac{1}{x-a}
\]
are partial: they have poles when the denominator vanishes.

The $\ph$-native setting changes the problem.  In the companion paper, the scalar grammar generated by
\[
    \B(x,y)=xy-y
\]
and the leaf
\[
    \ph=\frac{1+\sqrt5}{2}
\]
was shown to have exact constants
\[
    \Z[\ph]=\{a+b\ph:a,b\in\Z\}.
\]
Because $\ph$ is irrational, $\Z[\ph]$ is dense in the selected real embedding of $\R$; under the full Minkowski embedding into $\R^2$, it is discrete.  Hence one-coordinate density is achieved without division.  This removes one of the ordinary reasons for adding division early.

The guiding question of this note is therefore:
\begin{quote}
Once density no longer requires division, what remains of division?
\end{quote}

The proposed answer is that division splits into several $\ph$-native mechanisms:
\begin{center}
\begin{tabular}{p{0.27\linewidth}p{0.20\linewidth}p{0.42\linewidth}}
\toprule
Classical role & Classical notation & $\ph$-native replacement \\
\midrule
Unit scaling & $x/\ph$ & unit action $\B(\ph,x)$ \\
Fixed-denominator rate & $x/n$ & unevaluated localization class $[x:n]$ \\
Averaged value & $x/n$ as scalar & diagonal trace section $nq=x$ \\
Rational constants & $1/n$ & rate/localization layer or dense $\Z[\ph]$-approximants \\
Reciprocal functions & $1/x$ & local polynomial approximation \\
Quotient solving & $u/v$ & relation $u=qv$ \\
Derivative quotient & $\Delta f/\Delta x$ & factor projection \\
Normalization & $u/\|u\|$ & norm-one constraints \\
\bottomrule
\end{tabular}
\end{center}

The slogan is:
\begin{quote}
Division should not be imported wholesale into the $\ph$-native universe.  It should be decomposed into the roles it was performing.
\end{quote}

The resulting layers can be read as follows:
\begin{center}
\begin{tabular}{p{0.26\linewidth}p{0.62\linewidth}}
\toprule
Layer & Role \\
\midrule
Exact scalar layer & finite constants and polynomial expressions over $\Z[\ph]$ \\
Rate/localization layer & fixed-denominator classes such as $[x:n]\in \Z[\ph][1/n]$ \\
Analytic approximation layer & polynomial approximants on compact domains away from singularities \\
Relation-solving layer & equations such as $u=qv$ rather than primitive quotient terms \\
Factor-projection layer & derivative-like quotients as extraction of a known factor \\
\bottomrule
\end{tabular}
\end{center}

This note is deliberately not a full calculus paper.  It contains a worked golden-difference example, but the purpose of that example is to show how a derivative-like object can arise from factor projection rather than from a primitive quotient.  Full $\ph$-calculus, trigonometry, and analytic completion belong to later work.

\section{Background from the golden shadow algebra}

We recall only the pieces of the companion manuscript needed here.

\begin{definition}[Scalar core]
The scalar operation is
\[
    \B(x,y)=xy-y=y(x-1).
\]
The distinguished leaf is $\ph$, satisfying
\[
    \ph^2-\ph=1.
\]
\end{definition}

The basic scalar constructions are:
\[
    1=\B(\ph,\ph),
\]
\[
    0=\B(1,\ph),
\]
\[
    -x=\B(0,x),
\]
\[
    x-1=\B(x,1),
\]
\[
    x+1=\B\bigl(0,\B(\B(0,x),1)\bigr),
\]
and
\[
    xy=\B(x+1,y).
\]
Thus multiplication is scalar-native even though addition is not scalar-native.

The scalar-definable constants are exactly
\[
    \Z[\ph].
\]
The set $\Z[\ph]$ is countable, hence cannot equal $\R$, but it is dense in the selected real embedding of $\R$ because the fractional parts of $n\ph$ are dense modulo $1$ whenever $\ph$ is irrational.  This is a one-coordinate statement; under the full Minkowski embedding $a+b\ph\mapsto(a+b\ph,a+b\ps)$, $\Z[\ph]$ is discrete.  This is a standard consequence of irrational rotations/Kronecker-type density results.

The companion paper also introduces pair formation and the trace projection:
\[
    \Tr\langle x\mid y\rangle=x+y.
\]
This pair layer completes ordinary addition without making addition scalar-native.

For the present paper, the most important lesson is that the $\ph$-system starts from a dense, polynomial, division-free substrate.

\section{Division splits into roles}

Before replacing division, one must separate its roles.  In this paper we distinguish recurring roles: unit scaling, fixed-denominator rates, solved averages, exact constants, reciprocal functions, quotient solving, derivative quotients, and normalization.

\subsection{Unit scaling}

Some divisions merely multiply by an invertible constant:
\[
    \frac{x}{u}=u^{-1}x.
\]
When $u$ is a unit already present in the scalar ring, this is not a new operation.  In the golden ring, $\ph$ is a unit, and
\[
    \ph^{-1}=\ph-1.
\]

\subsection{Exact rational constants}

Division also introduces constants such as
\[
    \frac12,
    \qquad
    \frac13,
    \qquad
    \frac{a}{b}.
\]
These do not generally lie in $\Z[\ph]$.  For example,
\[
    \frac12\notin \Z[\ph].
\]
Thus exact rational scaling is not scalar-native.  There are several ways to respond: approximate such constants inside the dense ring \(\Z[\ph]\), carry the quotient as a rate class, or explicitly evaluate it in a controlled localization or averaging layer.  None of these is the same as a variable reciprocal.


\subsection{Rates, localization, and trace sections}\label{subsec:averaging}

Division by a fixed integer plays a different role from division by a variable expression.  The expression
\[
    x/n
\]
has no point at which the denominator can vanish as $x$ varies.  It is not a pole-producing reciprocal.  But it is also not generally an element of the integral scalar layer $R=\Z[\ph]$.

The clean way to record this is to treat fixed-denominator quotients first as rates.  For a multiplicative set $S\subset R\setminus\{0\}$, write
\[
    [x:s]_S,
    \qquad x\in R,
    \quad s\in S,
\]
for the rate represented by numerator $x$ and denominator $s$.  Two such representatives define the same rate when
\[
    [x:s]_S=[x':s']_S
    \quad\Longleftrightarrow\quad
    xs'=x's.
\]
Since $R=\Z[\ph]$ is an integral domain, this is the usual equality relation in the localization $S^{-1}R$.  Thus a rate is not a new scalar of $R$; it is an unevaluated localization class.  The denominator remains visible as dimensioned/count structure until a scalar coordinate is explicitly requested.

For a fixed integer $n\ge2$, the minimal denominator set is
\[
    S_n=\{1,n,n^2,\dots\},
\]
and the rate $[x:n]_{S_n}$ lives in the local layer
\[
    R[1/n].
\]
This is the rigorous version of carrying a quantity such as ``$5$ apples per $3$ people'': the object is represented by numerator-denominator data.  If one works in the broader integer-denominator localization, then
\[
    [5:3]=[10:6]=[15:9].
\]
If one works only in the minimal localization $S_3=\{1,3,9,\dots\}$, then $[5:3]=[15:9]$ is allowed while $[10:6]$ belongs to the broader localization.  The rate notation does not avoid localization; it delays evaluation by keeping the localization explicit.

The trace layer gives a complementary formulation when one wants a solved average.  For each integer $n\ge2$, define the diagonal map into an $n$-tuple sort by
\[
    \Delta_n(q)=\langle q\mid q\mid \cdots \mid q\rangle,
\]
and define the unnormalized $n$-trace by
\[
    \Tr_n\langle x_1\mid x_2\mid\cdots\mid x_n\rangle
    =x_1+x_2+\cdots+x_n.
\]
Then
\[
    \Tr_n(\Delta_n(q))=nq.
\]
Thus evaluating a rate as an average can be represented as a diagonal trace-section problem:
\[
    \operatorname{Avg}_n(x)=q
    \quad\Longleftrightarrow\quad
    \Tr_n(\Delta_n(q))=x.
\]
This is not a singular reciprocal.  It is a request for the scalar or module sort to be $n$-divisible.

\begin{proposition}[Rate localization and averaging completion]
Let $M$ be a $\Z[\ph]$-module.  Carrying fixed-denominator rates with denominators in $S_n=\{1,n,n^2,\dots\}$ corresponds to extension of scalars to
\[
    M[1/n]\cong M\otimes_{\Z[\ph]}\Z[\ph][1/n].
\]
For an integer denominator $n$, this agrees with the more elementary notation $M\otimes_{\Z}\Z[1/n]$.

A total operator $\operatorname{Avg}_n:M\to M$ satisfying
\[
    n\operatorname{Avg}_n(x)=x
\]
for all $x\in M$ is precisely a choice of section of multiplication by $n$.  Adjoining compatible sections for all $n\ge2$ is the rationalization step
\[
    M\longmapsto M\otimes_{\Z[\ph]}\Q(\ph),
\]
or equivalently $M\otimes_{\Z}\Q$ when only the underlying integer denominators are being inverted.
\end{proposition}

\begin{proof}
The equivalence relation above is the usual construction of the localization $S_n^{-1}M=M[1/n]$ as a $\Z[\ph][1/n]$-module.  The diagonal trace identity shows that a section of $\Tr_n\circ\Delta_n$ is the same thing as a section of multiplication by $n$.  Adjoining all rational averages is exactly extension of scalars from $\Z[\ph]$ to $\Q(\ph)$, equivalently from $\Z$ to $\Q$ on the underlying integer-denominator layer.
\end{proof}

The key distinction is therefore:
\[
    \frac{x}{n}
    \quad\text{is fixed-denominator rate/averaging, while}\quad
    \frac{1}{f(x)}
    \quad\text{is a variable reciprocal.}
\]
The first introduces a controlled localization or divisible module layer and remains total.  The second introduces genuine poles or domain restrictions.

\subsection{Reciprocal functions}

Division defines the function
\[
    x\mapsto \frac1x.
\]
This is a different role from constant scaling.  It introduces a pole at $x=0$.

\subsection{Quotient solving}

Division solves equations of the form
\[
    u=qv.
\]
The quotient $q=u/v$ is the solution, when it exists and is unique.

\subsection{Slope and normalization}

Division also appears in slopes,
\[
    \frac{\Delta f}{\Delta x},
\]
and in normalization,
\[
    u\mapsto \frac{u}{\|u\|}.
\]
These uses are conceptually distinct: slopes extract factors from differences, while normalization enforces a constraint such as norm one.

\begin{remark}
This paper does not reject division.  It argues that the $\ph$-native setting makes the overloading visible.  Once density is already present without division, each remaining role can be studied separately.
\end{remark}

\section{Unit actions are not general division}

The first simplification is that $x/\ph$ should not be treated as ordinary division.

\begin{proposition}[Golden contraction is scalar-native]
Let
\[
    C_\ph(x)=\B(\ph,x).
\]
Then
\[
    C_\ph(x)=(\ph-1)x=\frac{x}{\ph}.
\]
Thus contraction by $\ph^{-1}$ is a scalar polynomial operation.
\end{proposition}

\begin{proof}
By definition,
\[
    \B(\ph,x)=x(\ph-1).
\]
Since $\ph^2-\ph=1$, multiplying both sides by $\ph^{-1}$ gives
\[
    \ph-1=\ph^{-1}.
\]
Therefore
\[
    \B(\ph,x)=\ph^{-1}x.
\]
\end{proof}

So the notation
\[
    x/\ph
\]
should be avoided in the formal $\ph$-native presentation.  It should be replaced by
\[
    C_\ph(x)=\B(\ph,x).
\]
This is not a pole-producing reciprocal.  It is multiplication by a definable unit.

More generally, if $u$ is a unit in the scalar ring and $u^{-1}$ is definable, then multiplication by $u^{-1}$ is a unit action, not a new division primitive.

\section{Exact reciprocals and dense approximation}

Unit inverses are available when the unit is already present.  Non-unit reciprocals are different.  The constant $1/2$ is not in $\Z[\ph]$, so it is not a finite scalar constant of the grammar.

However, density supplies approximants.

\begin{proposition}[Approximating rational constants]
For every real number $r$ and every $\varepsilon>0$, there exists $c\in\Z[\ph]$ such that
\[
    |c-r|<\varepsilon.
\]
In particular, every rational constant can be approximated arbitrarily well by scalar-definable constants.
\end{proposition}

\begin{proof}
This is the density of $\Z[\ph]$ in the chosen real embedding of $\R$.  Since $\ph$ is irrational, the fractional parts of $b\ph$ are dense modulo $1$.  For a given $r$, choose $b\in\Z$ so that $b\ph$ is close to $r$ modulo $1$, and then choose $a\in\Z$ to correct the integer part.  Then $a+b\ph$ is as close to $r$ as desired.
\end{proof}

This gives approximate scaling without division.  But it is not free.  The integers $a,b$ needed to approximate a number to high precision may be large, and the expression
\[
    a+b\ph
\]
may involve cancellation between large terms.  Thus the $\ph$-native replacement avoids denominator singularities, but it may incur a syntactic and numerical complexity cost.

\begin{remark}[The syntactic complexity tax]
Density without division does not imply cheap representation.  Approximating $1/n$ or a general real number by $a+b\ph$ may require large coefficients.  This is not a singularity problem, but it is a size and conditioning problem.
\end{remark}

\section{Reciprocal functions as local approximation problems}

The reciprocal function
\[
    x\mapsto \frac1x
\]
has a genuine singularity at $x=0$.  A division-free grammar should not pretend to produce this function as a global polynomial.  No global polynomial replacement for $1/x$ is being claimed here.

The $\ph$-native replacement is local approximation.

\begin{proposition}[Golden-integer coefficient polynomial approximation]
Let $0<a<b$ and let $\varepsilon>0$.  There exists a polynomial
\[
    P(x)\in \Z[\ph][x]
\]
such that
\[
    \sup_{x\in[a,b]}\left|P(x)-\frac1x\right|<\varepsilon.
\]
\end{proposition}

\begin{proof}
By the Weierstrass approximation theorem, there is a real-coefficient polynomial
\[
    Q(x)=\sum_{j=0}^N r_jx^j
\]
uniformly approximating $1/x$ on $[a,b]$ within $\varepsilon/2$.  We now perturb the coefficients into $\Z[\ph]$.  Since $[a,b]$ is compact, the monomials appearing in $Q$ are uniformly bounded.  Let
\[
    M=\max_{0\le j\le N}\sup_{x\in[a,b]} |x|^j.
\]
If $M=0$ the claim is trivial; otherwise, by density of $\Z[\ph]$ in $\R$, choose $c_j\in\Z[\ph]$ such that
\[
    |c_j-r_j|<\frac{\varepsilon}{2(N+1)M}
\]
for every $0\le j\le N$.  Then, for every $x\in[a,b]$,
\[
\begin{aligned}
    \left|\sum_{j=0}^N(c_j-r_j)x^j\right|
    &\le \sum_{j=0}^N |c_j-r_j|\,|x|^j \\
    &< (N+1)\frac{\varepsilon}{2(N+1)M}M
     =\frac{\varepsilon}{2}.
\end{aligned}
\]
Thus $P(x)=\sum_{j=0}^N c_jx^j$ differs from $Q$ by less than $\varepsilon/2$ uniformly on $[a,b]$, and therefore differs from $1/x$ by less than $\varepsilon$ uniformly on $[a,b]$.
\end{proof}

The classical reciprocal gives an exact formula with a pole.  The $\ph$-polynomial approximation gives total finite expressions together with a domain of validity.  The singularity is not hidden; it is moved from the expression into approximation metadata:
\[
    \text{valid on }[a,b]\text{ with }0<a<b.
\]

This is the first major tradeoff:
\begin{center}
\begin{tabular}{lll}
\toprule
Method & Advantage & Cost \\
\midrule
Ordinary reciprocal $1/x$ & exact & pole at $x=0$ \\
$\Z[\ph]$-polynomial approximant & total finite expression & local domain and coefficient growth \\
\bottomrule
\end{tabular}
\end{center}

\section{Quotients as relations}

The algebraic content of a quotient can be expressed without a division operator.

\begin{definition}[Quotient relation]
Define
\[
    \Quot(u,v;q)
\]
to mean
\[
    u=qv.
\]
\end{definition}

In a field, when $v\ne0$, the relation has the unique solution
\[
    q=\frac{u}{v}.
\]
In a general ring, uniqueness requires cancellation assumptions, such as multiplication by $v$ being injective.  The $\ph$-native grammar need not contain a primitive expression $u/v$; it may instead treat the quotient as a relation to be solved or certified.

This matters because many quotient-like constructions are better understood relationally.  A slope, for example, need not be born as a fraction.  It can be the object $m$ satisfying
\[
    \Delta f=m\Delta x.
\]
This perspective leads directly to factor projection.

\section{Factor projection}

The companion manuscript, \emph{A Golden Shadow Algebra}, showed that addition becomes natural when represented as a projection from a pair:
\[
    \Tr\langle x\mid y\rangle=x+y.
\]
Here the analogous claim is:
\begin{quote}
Derivative-like quotients become natural when represented as factor projections from a difference object.
\end{quote}

\begin{definition}[Factor projection]
Let $A$ be a commutative ring and let $d\in A$ be fixed.  Suppose multiplication by $d$ is injective on $A$; for example, $A$ may be an integral domain and $d\ne0$.  A factor projection along $d$, when defined, sends an element $h\in (d)$ to the unique element $g$ satisfying
\[
    h=gd.
\]
We denote this operation by
\[
    \Fact_d(h)=g.
\]
\end{definition}

This is not division as a global operation.  It is a partial projection from the principal ideal $(d)$ to its coefficient along $d$, and uniqueness depends on the non-zero-divisor condition above.  If the factorization is known structurally, the projection is singularity-free at the symbolic level.

The derivative quotient
\[
    \frac{\Delta f}{\Delta x}
\]
should therefore be reinterpreted as:
\[
    \Fact_{\Delta x}(\Delta f).
\]

The next section develops the golden version of this idea.  It is a worked example and a preview of future $\ph$-calculus, not the main body of the present paper.

\section{Golden difference as a worked example}

Let
\[
    q=\ph^{-1}=\ph-1.
\]
The golden contraction is the unit action
\[
    C_\ph(x)=\B(\ph,x)=qx.
\]
Define the golden difference of a polynomial $f\in\Z[\ph][x]$ by
\[
    \Delta_\ph f(x)=f(x)-f(C_\ph x)=f(x)-f(qx).
\]
In particular,
\[
    \Delta_\ph x=x-qx=(1-q)x=(2-\ph)x.
\]
Notice that $x=0$ is a fixed point of the contraction.  A pointwise quotient
\[
    \frac{\Delta_\ph f(x)}{\Delta_\ph x}
\]
would therefore run into a $0/0$ expression at the origin.  Factor projection avoids this: it defines the derivative-like object as a formal polynomial factor, not as pointwise division.  In particular, although $\Delta_\ph x$ vanishes when evaluated at $x=0$, the identity
\[
    \Delta_\ph f=(\partial_\ph f)\Delta_\ph x
\]
lives in the polynomial ring; the factor $\partial_\ph f$ is determined globally before any pointwise evaluation is made.

\begin{theorem}[Golden factor derivative]
For every polynomial $f\in\Z[\ph][x]$, there is a unique polynomial $\partial_\ph f\in\Z[\ph][x]$ such that
\[
    \Delta_\ph f=(\partial_\ph f)\,\Delta_\ph x.
\]
Thus
\[
    \partial_\ph f=\Fact_{\Delta_\ph x}(\Delta_\ph f)
\]
is a well-defined factor projection on polynomials.
\end{theorem}

\begin{proof}
It suffices to check monomials.  For $f(x)=x^n$ with $n\ge1$,
\[
    \Delta_\ph x^n=x^n-(qx)^n=(1-q^n)x^n.
\]
Since
\[
    1-q^n=(1-q)(1+q+q^2+\cdots+q^{n-1}),
\]
we have
\[
\begin{aligned}
    \Delta_\ph x^n
    &= (1-q)(1+q+\cdots+q^{n-1})x^n \\
    &= \bigl(1+q+\cdots+q^{n-1}\bigr)x^{n-1}\,(1-q)x.
\end{aligned}
\]
But
\[
    \Delta_\ph x=(1-q)x.
\]
Substituting this displacement factor gives
\[
    \Delta_\ph x^n
    =\bigl(1+q+\cdots+q^{n-1}\bigr)x^{n-1}\Delta_\ph x.
\]
Hence $\Delta_\ph x^n$ is divisible by $\Delta_\ph x$ with quotient in $\Z[\ph][x]$.  Linearity over coefficients gives the result for all polynomials.  Uniqueness follows because $\Z[\ph][x]$ is an integral domain and $\Delta_\ph x\ne0$.  This is algebraic uniqueness in the polynomial ring, so it is unaffected by the fact that the displacement evaluates to zero at the fixed point $x=0$.
\end{proof}

\begin{corollary}[Monomial rule]
For $n\ge1$,
\[
    \partial_\ph x^n=[n]_\ph x^{n-1},
\]
where
\[
    [n]_\ph=1+q+q^2+\cdots+q^{n-1}\in\Z[\ph].
\]
\end{corollary}

This is the $q$-integer rule with $q=\ph^{-1}$, but here it appears without introducing a quotient operator.  The coefficient
\[
    [n]_\ph
\]
is a finite sum of scalar-definable constants.

\begin{proposition}[Twisted product rule]
For $f,g\in\Z[\ph][x]$,
\[
    \partial_\ph(fg)(x)=\partial_\ph f(x)\,g(x)+f(qx)\,\partial_\ph g(x).
\]
\end{proposition}

\begin{proof}
Compute
\[
    f(x)g(x)-f(qx)g(qx)
\]
as
\[
    \bigl(f(x)-f(qx)\bigr)g(x)+f(qx)\bigl(g(x)-g(qx)\bigr).
\]
Using
\[
    f(x)-f(qx)=(\partial_\ph f)(x)\Delta_\ph x
\]
and
\[
    g(x)-g(qx)=(\partial_\ph g)(x)\Delta_\ph x,
\]
then extracting the common factor $\Delta_\ph x$ gives the result.
\end{proof}

\begin{remark}[Relation to $q$-calculus]
This construction is formally the Jackson $q$-difference operator at $q=\ph^{-1}$, a classical object in $q$-calculus~\cite{Jackson1909,KacCheung2002}.  The claim here is not historical novelty of the operator.  The point is structural: the $q$-difference operator arises naturally when one rebuilds derivative-like behavior from a $\ph$-native unit action and factor projection, while refusing to introduce a primitive singular quotient.
\end{remark}

\section{Exponential preview: formal path before analytic completion}

The ordinary exponential is characterized by
\[
    \frac{d}{dx}e^x=e^x.
\]
A golden exponential should therefore be an eigenfunction of $\partial_\ph$:
\[
    \partial_\ph E_\ph=E_\ph.
\]

One might write a formal series
\[
    E_\ph(x)=\sum_{n\ge0}\frac{x^n}{[n]_\ph!},
\]
where
\[
    [n]_\ph!=[n]_\ph[n-1]_\ph\cdots[1]_\ph.
\]
But this notation reintroduces reciprocals into the coefficients.  To preserve the division-free philosophy, the formal path should come first.

Let
\[
    E_\ph(x)=\sum_{n\ge0}a_nx^n.
\]
The eigenfunction equation implies the recurrence
\[
    a_{n+1}[n+1]_\ph=a_n.
\]
This recurrence can be studied as a relation before solving it by explicit division.  It is therefore not yet a construction inside the finite scalar grammar.  Passing to an analytic completion, where reciprocal coefficients are allowed, is a later step.

Thus the exponential analogue is a preview of the broader program:
\begin{quote}
finite expression trees remain polynomial and total; analytic functions enter as completions or solutions of relations.
\end{quote}

\section{Normalization by constraint}

Division also appears in normalization:
\[
    u\mapsto \frac{u}{\|u\|}.
\]
A $\ph$-native alternative is to treat normalization as a constraint rather than an operation.

In the shadow algebra, norm-one objects satisfy
\[
    \Nm(u)=1.
\]
For the Gaussian pair
\[
    I=\langle i\mid -i\rangle,
\]
we have
\[
    \Tr(I)=0,
    \qquad
    \Nm(I)=1.
\]
This is the algebraic skeleton of rotation: trace zero and norm one.  Instead of producing a normalized object by dividing by its length, one studies the norm-one locus directly.

This prepares the next stage of the program.  Trigonometry may be better approached through norm-one shadow orbits than through primitive sine and cosine functions.

\section{Comparison with ordinary division}

The following table summarizes the proposed split.

\begin{center}
\begin{tabular}{p{0.27\linewidth}p{0.20\linewidth}p{0.42\linewidth}}
\toprule
Classical role & Classical form & $\ph$-native replacement \\
\midrule
Unit inverse & $x/\ph$ & $\B(\ph,x)$ \\
Fixed rate & $x/n$ & localization class $[x:n]$ \\
Solved average & $x/n$ as scalar & diagonal trace section $\Tr_n\Delta_n(q)=x$ \\
Rational constants & $1/n$ & rate/localization layer or $\Z[\ph]$-approximation \\
Reciprocal function & $1/x$ & local polynomial approximation \\
Quotient & $u/v$ & relation $u=qv$ \\
Derivative quotient & $\Delta f/\Delta x$ & factor projection \\
Normalization & $u/\|u\|$ & norm-one constraint \\
\bottomrule
\end{tabular}
\end{center}

This is the main conceptual claim of the paper.  In a $1$-leaf system, division is often needed early to escape the discreteness of $\Z$.  In the $\ph$-leaf system, density is already present in a chosen real coordinate.  Division is no longer the gateway to approximation; it becomes a collection of roles to be rebuilt more carefully.

\section{Open questions}

\begin{question}[Canonical golden-integer approximants]
Can one define a canonical approximation scheme
\[
    a_n+b_n\ph\to r
\]
for every real $r$?  Continued fractions, nearest-lattice methods, and modular rotations all suggest possible answers.
\end{question}

\begin{question}[Approximation cost]
How large must $a,b$ be to approximate $1/n$ or a general real constant within $\varepsilon$ by $a+b\ph$?  What is the correct complexity measure for $\ph$-integer approximants?
\end{question}

\begin{question}[Reciprocal approximation]
What are efficient $\Z[\ph]$-coefficient polynomial approximants to $1/x$ on compact domains away from zero?  How does coefficient growth compare with rational or Chebyshev approximations?
\end{question}

\begin{question}[Averaging sections]
Can rate classes $[x:s]_S$ and diagonal trace-section operators $\operatorname{Avg}_n$ be organized as a coherent typed localization/averaging layer compatible with pair trace, factor projection, and later phase normalization?  What is the right syntax for keeping rates unevaluated until scalar coordinates are explicitly requested?
\end{question}


\begin{question}[Factor projection as structure]
Can factor projection be formalized categorically in the same spirit as trace projection?  More generally, when should a classical quotient be replaced by a projection from a structured object?
\end{question}

\begin{question}[Golden exponentials]
What is the correct $\ph$-native exponential: a Jackson $q$-exponential at $q=\ph^{-1}$, a Fibonacci/Lucas trace exponential, or a family of related completions?  Which version best preserves the total, singularity-free finite expression layer?
\end{question}

\begin{question}[Toward shadow trigonometry]
Can trigonometric functions be reconstructed from norm-one shadow orbits?  Is the Gaussian pair $\langle i\mid -i\rangle$ the right entry point for a $\ph$-native trigonometric layer?
\end{question}

\section{Conclusion: division splits}

The $\ph$-native scalar grammar reaches a surprising starting point: exact finite constants are only $\Z[\ph]$, but this ring is dense in a chosen real embedding of $\R$, and all finite scalar expressions remain polynomial and total.  Thus density is achieved without division and without denominator singularities.

This changes the role of division.  Division is no longer a monolithic primitive required to reach a dense arithmetic substrate.  It splits into distinct mechanisms:
\[
\begin{gathered}
    \text{unit action},\quad \text{rate localization},\quad \text{averaging},\quad \text{approximation},\\
    \text{relation solving},\quad \text{factor projection},\quad \text{norm constraint}.
\end{gathered}
\]
Fixed integer division such as \(x/n\) belongs first to a rate/localization layer: it is represented by an unevaluated class \([x:n]\), and when a solved scalar average is required it becomes a diagonal trace section.  This enlarges the scalar or module sort, but it does not introduce variable-denominator singularities.

The golden contraction
\[
    C_\ph(x)=\B(\ph,x)
\]
is not division but unit action.  The golden derivative is not born from a quotient but from extracting the displacement factor
\[
    \Delta_\ph f=(\partial_\ph f)\Delta_\ph x.
\]

The resulting picture is not a rejection of classical division, analysis, or $q$-calculus.  It is a reordering of foundations.  In the $\ph$-native universe, finite expressions first form a one-coordinate dense, singularity-free algebraic substrate.  Classical analytic objects then enter later as completions, relations, or coordinatizations of structures already present.

\begin{thebibliography}{9}

\bibitem{SantosGoldenShadow}
V.~Santos.
\newblock \emph{A Golden Shadow Algebra: From a $\varphi$-Leaf Scalar Core to Conjugate-Pair Arithmetic}.
\newblock Exploratory manuscript, 2026.

\bibitem{Jackson1909}
F.~H. Jackson.
\newblock On $q$-functions and a certain difference operator.
\newblock \emph{Transactions of the Royal Society of Edinburgh}, 46(2):253--281, 1909.

\bibitem{KacCheung2002}
V.~Kac and P.~Cheung.
\newblock \emph{Quantum Calculus}.
\newblock Springer, 2002.

\bibitem{HardyWright2008}
G.~H. Hardy and E.~M. Wright.
\newblock \emph{An Introduction to the Theory of Numbers}.
\newblock 6th edition, Oxford University Press, 2008.

\bibitem{Apostol1974}
T.~M. Apostol.
\newblock \emph{Mathematical Analysis}.
\newblock 2nd edition, Addison-Wesley, 1974.

\bibitem{Lorentz1986}
G.~G. Lorentz.
\newblock \emph{Approximation of Functions}.
\newblock Chelsea Publishing, 1986.

\bibitem{ConradTraceNorm}
K.~Conrad.
\newblock \emph{Trace and Norm}.
\newblock Expository notes, University of Connecticut. Available at \url{https://kconrad.math.uconn.edu/blurbs/galoistheory/tracenorm.pdf}.

\end{thebibliography}

\end{document}
